\section{Programming languages: an overview}
\label{sec:overview}

\subsection{What is a programming language}

The course starts from a short definition.

\begin{definition}
A \textbf{programming language} is a system to write computer programs.
\end{definition}

The definition is short on purpose, and the rest of the course fills it in.
Figure~\ref{fig:pl-components} shows the parts that the first lecture adds to
it. The syntax of a language says how programs look: which symbols can be
used and how they combine into structures. The semantics says how programs
work: what type each expression has, which environment gives values to the
names, how names are bound, and how the program is represented internally
when it runs. Around both sits the tooling (editors, compilers, debuggers),
which does not change what a program means but decides how pleasant the
language is to use.

\begin{figure}[h]
\centering
\begin{forest}
  for tree={font=\small, s sep=4mm, l sep=6mm, align=center, draw, rounded corners=1pt, inner sep=3pt}
  [Programming language
    [{Syntax\\\footnotesize how programs look}
      [Symbols] [Structure]]
    [{Semantics\\\footnotesize how programs work}
      [Type system] [Environment] [Bindings] [{Internal\\representation}]]
    [Tooling]
  ]
\end{forest}
\caption{The components of a programming language, as presented in the first lecture.}
\label{fig:pl-components}
\end{figure}

Each block of these notes takes one of the boxes. Symbols and structure are
the subject of this block and the next one (grammars, automata, lexing,
parsing). Type systems, environments and bindings come after that, and
internal representation is the subject of the runtime block.

\subsection{What is a program}

Take this text from the lecture:
\begin{lstlisting}
fun arithmetic_logic(x,y) {
  a = 3+foo;
  x and y;
}
\end{lstlisting}
Whether it is a program depends on the language. A string is a program of a
language $L$ only if it follows the syntax of $L$, and it is a meaningful
program only if it also follows the semantic rules of $L$. Without fixing $L$
the question has no answer.

In the lecture the text is given to a parser for a small language whose
statements are assignments and expressions. Figure~\ref{fig:arith-logic} shows
the tree the parser builds. Function definition, parameters and the second
statement \texttt{x and y;} all fit the grammar. The first statement does not:
the grammar of that language has no \texttt{+} operator, so after
\texttt{a = 3} the parser finds \texttt{+foo} where it expects \texttt{;}. It
reports the error at \texttt{foo} (dashed box in the figure), skips the
offending tokens and continues. So the text is not a program of that
language, even though every token in it is valid.

\begin{figure}[h]
\centering
\begin{forest}
  arbol,
  [program
    [def
      [arithmetic\_logic] [(] [x] [{,}] [y] [)] [\{]
      [stat
        [a] [=] [expr [3]] [foo, draw, dashed, inner sep=2pt] [;]]
      [stat
        [expr [expr [x]] [and] [expr [y]]] [;]]
      [\}]
    ]
    [<EOF>]
  ]
\end{forest}
\caption{Parse tree for \texttt{arithmetic\_logic}. The parser cannot place
\texttt{foo} and marks it as an error.}
\label{fig:arith-logic}
\end{figure}

Passing the parser is not enough either. If the grammar had \texttt{+}, the
statement \texttt{a = 3+foo} would still need \texttt{foo} to be bound to
something, and \texttt{x and y} would need \texttt{x} and \texttt{y} to be
booleans. Those checks belong to the semantics, and different languages make
them at different moments.

The lecture shows this with a single line typed into seven interactive
interpreters (Table~\ref{tab:abc}). The text \texttt{a = b + c} is the same in
all of them, and none of them runs it.

\begin{table}[h]
\centering
\small
\begin{tabular}{l l l}
  \toprule
  Language & Response to \texttt{a = b + c} & When the problem is found \\
  \midrule
  JavaScript & \texttt{ReferenceError: b is not defined} & while running \\
  Ruby       & \texttt{NameError: undefined local variable or method `b'} & while running \\
  Python     & \texttt{NameError: name 'b' is not defined} & while running \\
  Elixir     & \texttt{undefined variable "b"}, \texttt{CompileError} & before running \\
  Haskell    & \texttt{Variable not in scope: b}, and the same for \texttt{c} & before running \\
  C++        & \texttt{Use of undeclared identifier 'a'} & before running \\
  Prolog     & \texttt{false.} & while running \\
  \bottomrule
\end{tabular}
\caption{The same line in seven languages. Each one reads it with its own semantics.}
\label{tab:abc}
\end{table}

The differences are not in the text but in the rules each language applies to
it. JavaScript, Ruby and Python look up \texttt{b} when the line is executed
and fail there. Elixir, Haskell and C++ check names before the program runs.
C++ complains about \texttt{a} first because it requires every variable to be
declared before its first use, including the one being assigned. Prolog does
not read \texttt{=} as assignment at all: it reads \texttt{a = b + c} as the
question ``does the atom \texttt{a} unify with the term \texttt{b + c}?'', and
the answer is no.

\subsection{Languages seen in class}

In the first lecture each student writes a function that computes the
factorial of a number, and then the same function is shown in several
languages. Table~\ref{tab:languages} summarizes them. The point is that the
same algorithm can look very different depending on what the language was
designed for.

\begin{table}[h]
\centering
\small
\begin{tabular}{l p{0.68\textwidth}}
  \toprule
  Language & What it shows \\
  \midrule
  Ada & Structured, statically typed, imperative and object-oriented. Strong
        typing, design by contract and concurrency are built in; the
        compiler is expected to catch errors before the program runs. The
        factorial uses subtypes \texttt{Positive} and \texttt{Natural}, so a
        negative argument is a type error. \\
  Prolog & Logic programming. A program is a set of facts and rules that
           define relations, and computation starts with a query. The
           factorial is two clauses relating \texttt{X} and \texttt{F}. \\
  Scratch & Block-based visual language for children aged 8 to 16. Programs
            are built by snapping blocks together, so syntax errors cannot
            be written. \\
  ASCII Dots & Esoteric language based on ASCII art. Dots travel along paths
               drawn with characters and are transformed by the operators
               they cross. \\
  CLU & Introduced abstract data types, call by sharing, iterators, multiple
        return values and type-safe parameterized types. Its factorial
        \texttt{signals negative} instead of returning a wrong value. \\
  UML & General-purpose visual modeling language with structure, behavior
        and interaction diagrams. \\
  Action! & Procedural language for Atari 8-bit computers, similar to
            ALGOL, compiled to fast 6502 code. Its factorial reports errors
            through a pointer argument because the result type cannot hold
            them. \\
  0815 & Esoteric language based on a queue; the factorial is written with
         jumps and labels such as \texttt{\}:f:} and \texttt{\^{}:f:}. \\
  Piet & Stack-based esoteric language whose programs are images. It uses 20
         colors; transitions between colors encode the operations. \\
  SCuLPT & Simple Cubic Language for Programming Tasks. Programs are built
           from tangible 3D pieces instead of typed text, and the language
           is Turing complete. \\
  \bottomrule
\end{tabular}
\caption{Languages used in the first lecture.}
\label{tab:languages}
\end{table}

For reference, here is the factorial in the two languages used for code in
these notes. In Scheme:
\codelines{Scheme}{scheme/factorial.scm}{2}{5}
and in Rascal:
\codelines{Rascal}{rascal/src/Factorial.rsc}{3}{3}
Both are recursive definitions by cases, and both files include tests.

\subsection{What makes a good programming language}

Pratt and Zelkowitz~\cite{pratt2001design} give eight criteria, and the
lecture follows them.

\paragraph{Clarity, simplicity and unity of concepts.} A language is a tool
for thinking about algorithms, writing them down and communicating them. It
should offer a small set of concepts, with simple and regular rules to combine
them, and it should be unambiguous. The dangling \texttt{else} of
Section~\ref{sec:cfl} is an example of a construct where the grammar alone
does not say what a program means.

\paragraph{Clarity of program structure.} The syntax should not mislead, and
the structure of a program should reflect the logical structure of the
algorithm.

\paragraph{Naturalness for the application.} The language should provide the
data structures, operations, control structures and syntax that the problem
needs. The date language of Section~\ref{sec:regular} is natural for dates and
useless for anything else, which is what a domain-specific language aims for.

\paragraph{Ease of extension.} Extensions should use simple mechanisms built
from primitive features of the language, so that they feel like part of it.

\paragraph{External support.} Tools for testing, editing and storing
programs, and good documentation, decide much of a language's success.

\paragraph{Ease of verification.} Programs can be proven correct with a formal
method, validated by tests that compare outputs against the specification, or
desk-checked by reading them. In all three cases, simple syntax and semantics
make the work easier.

\paragraph{Portability.} Programs should move between computing environments
without changes.

\paragraph{Efficiency and cost.} There are several costs: executing the
program (compiler optimizations and runtime support), translating it (fast
compilation matters for programs compiled often and run rarely), creating it
(design, coding, testing, modifying) and maintaining it.

MacLennan~\cite{maclennan1999principles} gives a longer list of design
principles, shown in Table~\ref{tab:maclennan}. Several of them come back
later in the course: orthogonality and regularity when we compare grammars,
information hiding when we discuss scopes, security when we discuss type
systems.

\begin{table}[h]
\centering
\small
\begin{tabular}{l p{0.7\textwidth}}
  \toprule
  Principle & Meaning \\
  \midrule
  Abstraction & Do not require something to be stated more than once; factor out recurring patterns. \\
  Automation & Automate mechanical, tedious or error-prone activities. \\
  Defense in depth & Have a series of defenses, so an error missed by one is caught by another. \\
  Information hiding & Users of a module get all the information they need and nothing more; implementers get all they need to implement it. \\
  Labeling & Do not require the user to know absolute positions in a list; give items meaningful labels. \\
  Localized cost & Users pay only for what they use. \\
  Manifest interface & All interfaces are apparent in the syntax. \\
  Orthogonality & Independent functions are controlled by independent mechanisms. \\
  Portability & Avoid features that depend on a particular machine. \\
  Preservation of information & Let the user state what they know and the compiler might need. \\
  Regularity & Regular rules without exceptions are easier to learn, use, describe and implement. \\
  Security & No program that violates the definition of the language should escape detection. \\
  Simplicity & A minimum number of concepts with simple rules for combining them. \\
  Structure & The static structure of the program corresponds to the dynamic structure of the computation. \\
  Syntactic consistency & Similar things look similar; different things look different. \\
  Zero-one-infinity & The only reasonable numbers are zero, one and infinity. \\
  \bottomrule
\end{tabular}
\caption{MacLennan's principles of language design.}
\label{tab:maclennan}
\end{table}

\subsection{A brief history}

The history of programming languages explains why some ideas exist. Three
periods are enough for this course.

The first languages were machine code and assembly (1940s and 1950s). A
program was a sequence of numeric opcodes, later written with mnemonic names.
One line was one machine instruction, and a program only ran on the processor
it was written for.

The early high-level languages came next (1950s and 1960s). FORTRAN (1957),
Lisp (1958) and COBOL (1959) separated source code from the hardware. FORTRAN
gave scientists expressions and arrays instead of registers. Lisp introduced
code as data, recursion as the main control structure and garbage collection.

From the 1970s on, Pascal and C brought structured control flow without
\texttt{goto}, Smalltalk and then C++ and Java brought objects, and ML and
Haskell brought typed functional programming with type inference. Each step
gave up some low-level control in exchange for a guarantee or for less work
for the programmer. Types and semantics, which come later in the course, are
the tools that state precisely what was gained.

\subsection{Design goals in tension}

The criteria of the previous section pull in different directions, and every
language resolves the conflicts in its own way.

\paragraph{Expressiveness and simplicity.} A more expressive language lets
programmers say more, often in fewer lines, but a larger language is harder to
learn, to implement and to reason about. This is the tension between the first
criterion (simplicity) and the third (naturalness for the application).

\paragraph{Safety and flexibility.} A safe language rejects some programs on
purpose. A type checker that rejects a program before it runs removes a class
of runtime errors, and it also rejects some programs that would have run
correctly. Table~\ref{tab:abc} shows both choices: Haskell rejects
\texttt{a = b + c} before running, Python runs it and fails.

\paragraph{Performance and abstraction.} Garbage collection, dynamic dispatch
and generic containers make programs easier to write and reuse, and they cost
time and memory. A Python list of integers holds pointers to boxed objects in
the heap; a C array of \texttt{int} is a contiguous block of machine integers.
The runtime block of the course returns to this difference.

\paragraph{Readability and writability.} A notation that is quick to write can
be slow to read. Verbose languages are slow to write and easy to read cold.
Most general-purpose languages lean toward readability, since code is read
more often than it is written. Domain-specific languages sometimes go the other
way, because their users already know the notation of the domain.

\subsection{Paradigms}

A \textbf{paradigm} is a way of thinking about computation: what a step is,
what state means, and how control flow is expressed. Most languages mix
paradigms.

In the imperative paradigm a program is a sequence of statements that change
the state; the basic unit is the assignment and control flow is explicit.
In the functional paradigm a program is an expression built from function
applications, and computation is evaluation; in the pure case there is no
mutable state. In the logic paradigm a program is a set of facts and rules and
running it means asking a query; Prolog in Table~\ref{tab:languages} is the
example. In the object-oriented paradigm a program is a set of objects that
group state and behavior and communicate by calling methods.

The same small problem, adding the even numbers of a list, looks different in
each style. Imperative, in Python:
\codelines{Python}{python/sum_evens.py}{1}{6}
The accumulator and the loop are explicit. Functional, in Haskell:
\codelines{Haskell}{haskell/SumEvens.hs}{1}{2}
There is no accumulator and no loop: \texttt{filter even} selects the
elements and \texttt{sum} adds them, and the two functions are composed.
Declarative, in SQL:
\codelines{SQL}{sql/sum_evens.sql}{3}{3}
The query states the result and leaves the strategy to the database engine.
A logic version in Prolog would be three clauses: the sum of the empty list is
0, an even head is added to the sum of the tail, and an odd head is skipped.

\subsection{Why formal definitions}

Everything in this section so far was argued in prose. Prose is enough to
start, but it does not settle disagreements. Whether \texttt{arithmetic\_logic}
is a program, or whether \texttt{a = b + c} is an error, can only be decided
with a precise definition of the language. The course builds that definition
in layers. Grammars say which strings are programs (Section~\ref{sec:grammars}).
Automata give a model of computation that recognizes them and let us prove
facts such as ``no finite automaton recognizes balanced parentheses''
(Sections~\ref{sec:automata} and~\ref{sec:regular}). Type systems and
semantics, later in the course, state what programs mean and which ones are
rejected.

\subsection{Exercises}

\begin{exercise}
Write a function that computes the factorial of a number in a language you
know. Then answer: what does your function do with a negative argument, and
how would Ada, CLU and Action! handle that case according to
Table~\ref{tab:languages}?
\end{exercise}

\paragraph{Solution.} The recursive definition used in these notes returns 1
for every $n < 2$, so it returns 1 for negative numbers, which is not a
factorial. Ada declares the parameter as \texttt{Positive}, so a negative
argument is a type error. CLU declares \texttt{signals (negative)} and raises
that signal. Action! writes an error code through the pointer \texttt{err}
(1 for negative, 2 for a result that does not fit) and returns 0.

\begin{exercise}
Is \texttt{arithmetic\_logic} a program? Answer for the grammar of
Figure~\ref{fig:arith-logic}, and then for a grammar that also has
\texttt{+}. In the second case, what else would the language need to check?
\end{exercise}

\paragraph{Solution.} For the grammar of the figure it is not a program: the
first statement cannot be derived because there is no rule for \texttt{+}.
With \texttt{+} in the grammar the text is syntactically a program. The
language would still need to check that \texttt{foo} is bound (a scope rule),
that \texttt{3+foo} adds two numbers and that \texttt{x and y} combines two
booleans (type rules). Whether \texttt{a} needs a declaration depends on the
language, as the C++ row of Table~\ref{tab:abc} shows.

\begin{exercise}
For each language in Table~\ref{tab:abc}, say whether the problem with
\texttt{a = b + c} is found by the syntax or by the semantics, and explain why
Prolog answers \texttt{false} instead of reporting an error.
\end{exercise}

\paragraph{Solution.} In all seven languages the line is syntactically valid;
every response comes from the semantics. Six languages complain about a name
with no binding, at run time or before it. Prolog reads \texttt{=} as
unification. The left side is the atom \texttt{a} and the right side is the
compound term \texttt{+(b, c)}; an atom unifies only with itself or with a
variable, so unification fails and the query answers \texttt{false}.

\begin{exercise}
Choose two of Pratt and Zelkowitz's criteria and one language from
Table~\ref{tab:languages}. Explain how the language satisfies or fails each
criterion, with a concrete example.
\end{exercise}

\begin{exercise}
Pick two languages you know. For each of the four tensions of this section,
say which way each language leans, with an example from each.
\end{exercise}
